\problem{1}
\textbf{[Отбор 2025, 8.2]}
Точка $L$ - середина гипотенузы $A B$ прямоугольного треугольника $A B C$. На стороне $B C$ взята точка $N$ так, что $B N=4$ см., $N C=2$ см. Докажите, что $\angle B A N=\angle C L N$.

\problem{2}
\textbf{[Отбор 2023, 8.3]} В прямоугольнике $A B C D$ выбрана точка $E$ на стороне $A B$ и точка $F$ на отрезке $C E$ так, что треугольники $A D E$ и $C D F$ имеют равные площади. Докажите, что треугольники $B C E$ и $D E F$ также имеют равные площади.

\problem{3}
\textbf{[Отбор 2022, 8.3]}
В треугольнике $A B C$ на отрезке $B C$ отмечены точки $D$ и $E$ такие, что $B D= C E$ и $\angle B A D=\angle C A E$. Докажите, что треугольник $A B C$ равнобедренный.

\problem{4}
\textbf{[Отбор 2024, 8.4]}
В треугольнике $A B C$ с биссектрисой $C F$ выполняется соотношение $2 \angle C B A= 3 \angle A C B$. Точки $D$ и $E$ лежат на стороне $A C$, причём $B D$ и $B E$ делят $\angle C B A$ на три равных угла. Известно, что $D$ лежит между точками $A$ и $E$. Докажите, что $B E$ и $D F$ параллельны.

\problem{5}
\textbf{[Отбор 2021, 8.2]}
Дан параллелограмм $A B C D$. На стороне $A D$ выбрана точка $P$ так, что $B P$ является биссектрисой угла $B$. Найдите длины сторон параллелограмма $A B C D$, если $B P=C P=6$ и $P D=5$.

\problem{6}
\textbf{[Отбор 2021, 8.5]}
Квадрат со стороной 3 и прямоугольник с целыми сторонами, одна из которых на 2 длиннее другой, расположены так, что их стороны параллельны, а центры совпадают. Стороны квадрата продлили так, как показано на рисунке. Может ли площадь заштрихованной области быть чётным числом?
